\documentclass[aspectratio=169]{beamer}

\usetheme{Madrid}
\usecolortheme{default}
\setbeamertemplate{navigation symbols}{}

\usepackage{amsmath,amssymb,bm}
\usepackage{mathtools}

% -------------------- Notation --------------------
\newcommand{\x}{\bm{x}}
\newcommand{\uu}{\bm{u}}
\newcommand{\pp}{p}
\newcommand{\Rey}{\mathrm{Re}}
\newcommand{\avg}[1]{\overline{#1}}
\newcommand{\up}{\uu'}
\newcommand{\ubar}{\avg{\uu}}
\newcommand{\grad}{\nabla}
\newcommand{\Div}{\nabla\cdot}
\newcommand{\Lap}{\Delta}

\title[Turbulence Modeling]{Turbulence Modeling: RANS, LES, and Hybrid Ideas}
\author{(adapted from Steve Brunton)}
\date{}

\begin{document}

% ------------------------------------------------------------
\begin{frame}
  \titlepage
\end{frame}

% ------------------------------------------------------------
\begin{frame}{Why turbulence modeling? (design + HPC)}
Engineering design often needs \textbf{quantities of interest}:
\begin{itemize}
  \item lift and drag on wings / bodies,
  \item separation and stall prediction,
  \item turbine performance and noise,
  \item wake interactions (wind farms, vehicles).
\end{itemize}

\begin{block}{Core tension}
High fidelity $\Rightarrow$ expensive. \\
Design optimization/control requires \emph{many} simulation runs $\Rightarrow$ we need models that are cheaper than DNS.
\end{block}

\vspace{1em}
\textbf{Figure placeholder:} DNS movie / boundary layer turbulence visualization.
\end{frame}

% ------------------------------------------------------------
\begin{frame}{DNS is amazing... and usually too expensive}
Direct Numerical Simulation (DNS) resolves \emph{all} relevant space/time scales.

\begin{itemize}
  \item Even a \textbf{small} region (e.g., section of an airfoil) can require massive resources.
  \item A full aircraft / full automobile at realistic Reynolds numbers: typically out of reach.
\end{itemize}

\begin{block}{Cost scaling intuition}
As Reynolds number increases, the range of scales widens (largest eddies $\to$ smallest eddies), so grid/time-step requirements explode.
\end{block}

\textbf{Figure placeholder:} mesh + reported core-hours (e.g., ``35 million core hours'') from the example.
\end{frame}

% ------------------------------------------------------------
\begin{frame}{Goal of turbulence modeling}
Instead of resolving everything, we seek \textbf{approximations} that yield:
\begin{itemize}
  \item acceptable accuracy for engineering quantities,
  \item much lower computational cost,
  \item robustness across many parameter changes (geometry, speed, AoA, etc.).
\end{itemize}

\begin{block}{Two major families}
\begin{itemize}
  \item Reynolds-Averaged Navier--Stokes (RANS): model \emph{statistics} of turbulence.
  \item Large Eddy Simulation (LES): resolve large eddies, model subgrid scales (SGS).
\end{itemize}
\end{block}

\textbf{Preview:} Hybrid ideas (DES) mix RANS + LES.
\end{frame}

% ------------------------------------------------------------
\begin{frame}{Review: Reynolds decomposition (RANS entry point)}
Decompose the velocity (and pressure):
\[
\uu(\x,t)=\ubar(\x)+\up(\x,t),
\qquad
\pp(\x,t)=\avg{\pp}(\x)+\pp'(\x,t),
\qquad
\avg{\up}=0,\ \avg{\pp'}=0.
\]

Start from incompressible Navier--Stokes (nondimensional):
\begin{align*}
\partial_t \uu + (\uu\cdot\grad)\uu &= -\grad \pp + \frac{1}{\Rey}\Lap\uu,\\
\Div\uu &= 0.
\end{align*}

\begin{block}{From the continuity equation}
\[
\Div\ubar=0,\qquad \Div\up=0.
\]
\end{block}
\end{frame}

% ------------------------------------------------------------
\begin{frame}{Review: where closure appears (momentum)}
After averaging momentum you get mean-flow equations with extra terms:
\[
(\ubar\cdot\grad)\ubar
=
-\grad \avg{\pp}
+\frac{1}{\Rey}\Lap \ubar
-\Div \bm{R},
\qquad
\bm{R}=\avg{\up\,\up^T}.
\]

\begin{block}{Reynolds stresses}
\[
\bm{R}=
\begin{bmatrix}
\avg{u'u'} & \avg{u'v'} & \avg{u'w'}\\
\avg{v'u'} & \avg{v'v'} & \avg{v'w'}\\
\avg{w'u'} & \avg{w'v'} & \avg{w'w'}
\end{bmatrix}.
\]
\end{block}

\textbf{Closure problem:} $\bm{R}$ is unknown and depends on the turbulent fluctuations.
We need $\bm{R}\approx \mathcal{M}(\ubar,\grad\ubar,\ldots)$.
\end{frame}

% ------------------------------------------------------------
\begin{frame}{Why not write an exact equation for $\bm{R}$? (hierarchy)}
If you try to derive transport equations for the Reynolds stresses by multiplying the fluctuating equations and averaging, you get terms involving \textbf{triple correlations} (third moments):
\[
\avg{u'_i u'_j u'_k}, \quad \text{etc.}
\]

Then equations for third moments involve fourth moments, and so on.

\begin{block}{The never-ending ladder}
Second moments depend on third moments \\
Third moments depend on fourth moments \\
$\vdots$ \\
This is a genuine mathematical closure issue.
\end{block}

\textbf{Takeaway:} practical modeling requires assumptions/approximations.
\end{frame}

% ------------------------------------------------------------
\begin{frame}{Stress form in Einstein notation (compact RANS statement)}
A common compact form emphasizes ``stress divergence'' structure.
(Here, indices $i,j\in\{1,2,3\}$; repeated indices are summed.)

\[
\partial_t \avg{u_i} + \avg{u_j}\,\partial_j \avg{u_i}
=
-\partial_i \avg{p}
+\frac{1}{\Rey}\partial_j\partial_j \avg{u_i}
-\partial_j \avg{u_i' u_j'}.
\]

\begin{block}{Interpretation}
Pressure, viscosity, and Reynolds stresses all appear as stress/flux divergences.
\end{block}
\end{frame}

% ------------------------------------------------------------
\begin{frame}{Turbulent kinetic energy (TKE) as a modeling lens}
Define turbulent kinetic energy:
\[
k := \frac12\,\avg{u_i' u_i'} = \frac12\left(\avg{(u')^2}+\avg{(v')^2}+\avg{(w')^2}\right).
\]

TKE helps track the \textbf{energy exchange} between:
\begin{itemize}
  \item mean-flow kinetic energy,
  \item turbulent fluctuations.
\end{itemize}

\begin{block}{Conceptual idea}
Many closures aim to model $\bm{R}$ by modeling how $k$ is produced, transported, and dissipated.
\end{block}

\textbf{Figure placeholder:} energy budget / ``competing terms'' schematic.
\end{frame}

% ------------------------------------------------------------
\begin{frame}{TKE equation: anatomy (high-level)}
A typical TKE transport equation has the structure:
\[
\frac{\partial k}{\partial t} + \avg{u_j}\,\partial_j k
=
\underbrace{\mathcal{P}}_{\text{production}}
+\underbrace{\mathcal{T}}_{\text{transport}}
+\underbrace{\mathcal{D}}_{\text{diffusion}}
-\underbrace{\varepsilon}_{\text{dissipation}}.
\]

\begin{block}{Production term (key intuition)}
Production is driven by \textbf{mean shear}:
\[
\mathcal{P}\ \sim\ -\avg{u_i' u_j'}\,\partial_j \avg{u_i}.
\]
Large gradients in the mean flow (shear layers, separation) pump energy into turbulence.
\end{block}

\textbf{Use this in words:} ``Separation $\Rightarrow$ strong shear $\Rightarrow$ more turbulence production.''
\end{frame}

% ------------------------------------------------------------
\begin{frame}{Eddy viscosity modeling: simplest closure idea}
Old but still foundational idea: turbulence acts like an \textbf{effective viscosity}.

\begin{block}{Boussinesq hypothesis (eddy viscosity)}
Model Reynolds stresses via
\[
-\avg{u_i' u_j'}
\approx
2\nu_t S_{ij} - \frac{2}{3}k\,\delta_{ij},
\qquad
S_{ij}:=\frac12(\partial_i \avg{u_j}+\partial_j \avg{u_i}).
\]
\end{block}

\begin{itemize}
  \item $\nu_t$ is the turbulent (eddy) viscosity.
  \item The $-\tfrac{2}{3}k\,\delta_{ij}$ term corrects the isotropic part (normal stresses).
\end{itemize}

\textbf{Figure placeholder:} channel flow / boundary layer sketch + shear.
\end{frame}

% ------------------------------------------------------------
\begin{frame}{Eddy viscosity model 0: constant $\nu_t$}
Most basic choice:
\[
\nu_t = \text{constant}.
\]

Pros:
\begin{itemize}
  \item extremely simple,
  \item sometimes works acceptably after tuning for a specific configuration.
\end{itemize}

Cons:
\begin{itemize}
  \item not predictive across geometries/conditions,
  \item cannot capture strong spatial variation (near-wall, separation, etc.).
\end{itemize}
\end{frame}

% ------------------------------------------------------------
\begin{frame}{Prandtl mixing length model: spatially varying $\nu_t$}
Prandtl's correction: $\nu_t$ varies with local shear and a mixing length $\ell_m$.

A common form:
\[
\nu_t \sim \ell_m^2\,\left|\frac{\partial \avg{u}}{\partial y}\right|.
\]

\begin{itemize}
  \item Useful for wall-bounded flows (boundary layers, channels).
  \item Requires empirical/physical choices for $\ell_m$ (near-wall behavior matters).
\end{itemize}

\textbf{Figure placeholder:} boundary layer thickness growth and mixing length profile.
\end{frame}

% ------------------------------------------------------------
\begin{frame}{Two-equation models: $k$--$\varepsilon$ and $k$--$\omega$}
Modern ``workhorse'' RANS models solve extra PDEs.

\begin{block}{$k$--$\varepsilon$ idea}
Solve for:
\begin{itemize}
  \item $k$ (turbulent kinetic energy),
  \item $\varepsilon$ (dissipation rate of $k$).
\end{itemize}
Then set eddy viscosity via:
\[
\nu_t = C_\mu \frac{k^2}{\varepsilon}
\quad\text{(often with }C_\mu\approx 0.09\text{ empirically).}
\]
\end{block}

\textbf{Point to emphasize:} better than constant/mixing-length in many cases, but still imperfect.
\end{frame}

% ------------------------------------------------------------
\begin{frame}{Known weakness: adverse pressure gradients + separation}
Many standard RANS closures struggle with separated flows.

\begin{itemize}
  \item Adverse pressure gradient can cause flow reversal near the wall.
  \item Mean profile shows a \textbf{recirculation/separation bubble}.
\end{itemize}

\begin{block}{Why it matters}
Predicting separation and stall is critical for aircraft performance and safety.
\end{block}

\textbf{Figure placeholder:} separation bubble sketch + boundary layer velocity profiles (show reversal).
\end{frame}

% ------------------------------------------------------------
\begin{frame}{Spalart--Allmaras: practical aerospace workhorse}
Spalart--Allmaras (SA) is a widely used \textbf{one-equation} model, tuned for aerospace applications.

\begin{itemize}
  \item More efficient than two-equation models.
  \item Strong historical impact in CFD for lift/drag prediction.
\end{itemize}

\begin{block}{What to say out loud}
``SA is a compromise: relatively cheap, reasonably accurate for many aerodynamic flows, and widely validated in industry.''
\end{block}

\textbf{Figure placeholder:} ``turbulence model options'' screenshot / SA citation mention.
\end{frame}

% ------------------------------------------------------------
\begin{frame}{Transition: a different perspective (LES)}
So far: RANS models \textbf{all} turbulence statistically.

LES takes a different view:
\begin{itemize}
  \item Resolve \textbf{large} energetic eddies on the grid,
  \item Model only \textbf{subgrid} (small) scales.
\end{itemize}

\begin{block}{Why do this?}
Large eddies dominate transport and coherent structures (especially in separated flows),
so resolving them often increases accuracy.
\end{block}
\end{frame}

% ------------------------------------------------------------
\begin{frame}{LES concept: filtering / coarse-graining}
LES applies a spatial filter (low-pass) to the Navier--Stokes equations:
\[
\uu = \tilde{\uu} + \uu_{\mathrm{sgs}},
\]
where $\tilde{\uu}$ is the resolved (filtered) velocity.

\begin{itemize}
  \item Solve filtered equations on a coarser grid.
  \item Nonlinearity couples scales $\Rightarrow$ a closure term appears (SGS stress).
\end{itemize}

\begin{block}{Subgrid-scale (SGS) closure}
We must model the effect of unresolved scales on resolved scales (analogous to RANS closure, but ``below the filter'').
\end{block}

\textbf{Figure placeholder:} coarse grid boxes / ``under the hood'' schematic.
\end{frame}

% ------------------------------------------------------------
\begin{frame}{Kolmogorov cascade picture (why SGS looks like viscosity)}
Energy moves from large to small scales (heuristic cascade):
\begin{itemize}
  \item production at large scales (driven by mean shear / instabilities),
  \item inertial transfer across scales,
  \item dissipation at the smallest scales.
\end{itemize}

\begin{block}{LES modeling idea}
Do \emph{not} resolve deep into the dissipation range; treat unresolved small scales as an effective dissipation (often via eddy viscosity SGS models).
\end{block}

\textbf{Figure placeholder:} Kolmogorov energy spectrum / cascade diagram.
\end{frame}

% ------------------------------------------------------------
\begin{frame}{LES vs RANS: accuracy vs cost}
\begin{columns}[T,onlytextwidth]
\column{0.5\textwidth}
\textbf{RANS}
\begin{itemize}
  \item cheaper (often practical for design loops),
  \item solves mean field + closure PDEs,
  \item may struggle with strong separation / unsteady structures.
\end{itemize}

\column{0.5\textwidth}
\textbf{LES}
\begin{itemize}
  \item more detailed flow structures,
  \item typically better for separation / wakes,
  \item much more expensive (often needs HPC).
\end{itemize}
\end{columns}

\vspace{0.8em}
\textbf{Figure placeholder:} side-by-side LES (top) vs unsteady RANS (bottom) visualization for wake/separation.
\end{frame}

% ------------------------------------------------------------
\begin{frame}{Large-scale LES application example}
LES enables simulations at scales where DNS is impossible, e.g.:
\begin{itemize}
  \item wind-farm wake interactions across kilometers,
  \item atmospheric boundary layer transport,
  \item renewable energy control and optimization studies.
\end{itemize}

\begin{block}{How to frame it}
LES is not ``exact,'' but often captures the large-scale mixing and unsteady dynamics that matter for engineering decisions.
\end{block}

\textbf{Figure placeholder:} wind farm LES vorticity/wake visualization.
\end{frame}

% ------------------------------------------------------------
\begin{frame}{Hybrid idea: Detached Eddy Simulation (DES)}
DES attempts to combine strengths of RANS and LES:
\begin{itemize}
  \item use RANS near walls / attached regions (efficiency),
  \item switch to LES in massively separated wakes (accuracy).
\end{itemize}

\begin{block}{Mental model}
``RANS where the flow is comparatively steady/attached; LES where large unsteady eddies dominate.''
\end{block}

\textbf{Figure placeholder:} schematic showing RANS region near wall + LES region in wake.
\end{frame}

% ------------------------------------------------------------
\begin{frame}{Where machine learning enters (preview)}
Both RANS and LES hinge on closure:
\begin{itemize}
  \item RANS: model Reynolds stresses $\avg{u_i' u_j'}$,
  \item LES: model subgrid stresses / SGS fluxes.
\end{itemize}

\begin{block}{ML promise (high-level)}
Use high-fidelity data (DNS/experiments) to learn improved closures:
\[
\text{closure} \approx \text{function of resolved/mean features},
\]
potentially improving accuracy in hard regimes (separation, adverse pressure gradients).
\end{block}

\textbf{Next lecture (as in transcript):} deep learning for closure models.
\end{frame}

% ------------------------------------------------------------
\begin{frame}{Summary}
\begin{itemize}
  \item DNS is accurate but often too expensive for real design loops.
  \item RANS: average the equations $\Rightarrow$ Reynolds stresses appear $\Rightarrow$ closure problem.
  \item Eddy viscosity modeling (Boussinesq) is the simplest closure; mixing-length, $k$--$\varepsilon$, SA extend it.
  \item LES: filter the equations and resolve large eddies; model SGS effects.
  \item LES is often more accurate (especially for separated flows) but typically much more expensive.
  \item DES and other hybrids aim to mix RANS efficiency with LES accuracy.
\end{itemize}
\end{frame}

\end{document}
